\documentclass[11pt]{article}
\usepackage[margin=1in]{geometry}
\usepackage{booktabs}
\usepackage{amsmath}
\usepackage{url}
\usepackage{hyperref}
\usepackage{listings}
\usepackage{xcolor}
\lstset{basicstyle=\ttfamily\small, breaklines=true, frame=single, framesep=4pt}

\title{toolc: An Optimizing Compiler for Agent Tool Surfaces}
\author{Bryan Healey\\ \texttt{bryan@toolc.dev}}
\date{Draft, August 2026}

\begin{document}
\maketitle

\begin{abstract}
Language-model agents increasingly act through tools exposed via the Model
Context Protocol (MCP). In practice, agents connect to several servers at
once, and the naive federation of their catalogs scales poorly: every tool
definition occupies context, near-duplicate tools degrade selection accuracy,
and multi-step workflows must be improvised call-by-call. We present
\textsc{toolc}, an optimizing compiler that treats the federated tool surface
as a compilation target. \textsc{toolc} introspects downstream servers (or
synthesizes catalogs directly from OpenAPI specifications and human-readable
API documentation) into a capability-graph intermediate representation, then
applies a sequence of independently toggleable passes: dead-tool elimination,
LLM-assisted description rewriting with intra- and inter-server
disambiguation, consolidation of near-duplicate tool families into
deterministic facade tools, macro inlining, and an optional retrieval-based
selection surface. The compiled artifact is served as a single drop-in MCP
endpoint with full call logging, serve-time result compaction, and
error-driven re-optimization. In an evaluation over 137 tools from six live servers (60 tasks
$\times$ 3 trials $\times$ two agent models; 1{,}800 judged trials), the
compiled surface with consolidation matched the raw federation's success for
a frontier agent (96.1\% vs.\ 96.7\%, overlapping intervals) at $5\times$
lower cost per completed task, and \emph{improved} cross-source arbitration
(90\% $\rightarrow$ 100\% on overlap tasks); for a small agent the same
compression cost nine points, showing the benefit is capability-gated. Every optimization
is a measured, reversible configuration choice rather than a fixed
architecture.
\end{abstract}

\section{Introduction}

Tool use has become the dominant mechanism by which language-model agents act
on the world, and the Model Context Protocol has emerged as a standard
transport for exposing tools to agents. The unit of deployment, however, is
the \emph{server}, not the \emph{task}: a working agent setup typically
federates several servers, each contributing tens of tool definitions written
by different authors with different conventions.

This federation is an unoptimized program. Its costs are threefold:

\begin{enumerate}
\item \textbf{Context.} Every served tool definition is prompt input on every
turn. In our reference federation, 131 definitions consume roughly 24{,}000
tokens before the agent reads the user's request.
\item \textbf{Selection accuracy.} Catalogs contain near-duplicates
(\texttt{search\_x}, \texttt{find\_x}, \texttt{get\_x}), machine-generated
descriptions, and cross-server overlap (two news sources). The agent's tool
choice degrades with catalog entropy, and errors cost round trips.
\item \textbf{Orchestration.} Common multi-step workflows (find an entity,
fetch its records, extract a field) are re-derived by the agent on every
occurrence, at inference prices.
\end{enumerate}

Our thesis is that these are compiler problems. The tool surface admits an
intermediate representation, semantics-preserving transformations over that
representation, and a runtime that executes the compiled program while
observing it for further optimization. \textsc{toolc} realizes this thesis
end to end: it is available as an Apache-2.0 engine and as a hosted service
that adds managed endpoints, observability, and usage-based billing.

This paper makes four contributions:
\begin{enumerate}
\item A capability-graph IR for federated tool surfaces that preserves
original definitions as immutable facts and expresses every transformation as
a reviewable overlay (\S\ref{sec:ir}).
\item A pass architecture comprising dead-tool elimination, LLM-assisted
rewriting, deterministic consolidation of tool families into facades (within
and across servers), macro inlining, and a retrieval-based selection surface;
each pass is independently toggleable and its effect independently measurable
(\S\ref{sec:passes}).
\item A serving runtime with logged dispatch, serve-time result compaction,
and profile-guided re-optimization that converts recurring runtime errors into
proposed compiler fixes (\S\ref{sec:runtime}).
\item A judged, ablated evaluation methodology and pilot results showing that
compilation dominates the naive federation on both cost and accuracy axes
(\S\ref{sec:eval}).
\end{enumerate}

\section{Related Work}

\textbf{Tool learning and selection.} A substantial literature addresses
tool-augmented language models and tool selection at scale, including
retrieval over large tool inventories and hierarchical selection. Our
selection pass adopts retrieval (BM25 over compiled descriptions) as one
point on a measured cost/accuracy frontier rather than as the architecture.

\textbf{API-to-agent interfaces.} Automatic conversion of OpenAPI
specifications into agent tools is widely practiced and widely observed to
produce poor surfaces; our composition frontend embraces the mechanical
conversion precisely because the compiler owns making the result good, and
our evaluation includes a deliberately naive machine-generated server as a
worst-case input.

\textbf{Prompt and context optimization.} Description rewriting relates to
automatic prompt optimization; result compaction relates to context
summarization. \textsc{toolc} differs in binding these to a persistent,
versioned artifact with provenance (original text is never destroyed) and in
validating LLM proposals with deterministic checks before they take effect.

\textbf{Compiler framing.} We borrow the classical structure --- frontend,
IR, passes, emitter, profile-guided optimization --- and find it maps cleanly:
tool catalogs are modules, near-duplicate families are common subexpressions,
macros are inlined procedures, the selection surface is lazy loading, and
runtime error clustering is PGO.

\section{The Capability-Graph IR}\label{sec:ir}

A compiled surface is a graph $G = (S, T, E)$: sources $S$ (downstream
servers with catalog hashes), tools $T$, and data edges $E$ (produced-by /
consumed-by relations contributed by macros). Each tool node carries:

\begin{itemize}
\item identity (\texttt{source:name}), the \emph{original} description and
JSON-Schema input schema, immutable once introspected;
\item \emph{overlays}: the rewritten description, visibility
(\texttt{hidden} with a reason), and surface placement (\texttt{surfaced});
\item a kind: \texttt{passthrough}, \texttt{facade}, \texttt{macro}, or
\texttt{meta};
\item for facades, a deterministic routing table from action names to member
tool ids.
\end{itemize}

Two properties matter. \textbf{Provenance}: originals are never mutated, so
any transformation is auditable and reversible, and user-facing consoles can
display exact before/after diffs. \textbf{Content addressing}: the graph
serializes deterministically and is versioned by content hash, so a served
artifact is immutable and cache keys (e.g., for rewrites) are stable.

Frontends produce this IR from three inputs: live MCP introspection, OpenAPI
specifications (each operation becomes a passthrough tool with constraints
--- enums, bounds, defaults --- carried into the schema), and human-readable
documentation, from which a draft specification is synthesized by an LLM
under never-invent rules and then \emph{probe-validated}: parameterless GET
endpoints are called against the live API, and endpoints that return 404 are
flagged as probable hallucinations before anything compiles.

\section{Optimization Passes}\label{sec:passes}

Passes run in a fixed order, each mapping graph to graph; every pass emits a
structured diff for review. All passes are user-toggleable configuration.

\textbf{Dead-tool elimination} hides tools matched by user globs or by
built-in noise patterns (health checks, debug endpoints).

\textbf{Rewrite} regenerates descriptions with an LLM, batched per server
with full sibling context plus a digest of the rest of the pool, under
instructions to disambiguate against both siblings and overlapping tools on
other servers, to surface parameter constraints, and never to invent
capabilities. Proposals are cached by (tool, original-description hash,
prompt version) and can be gated on human review.

\textbf{Consolidate} merges families of near-duplicate tools into
\emph{facade} tools with an \texttt{action} discriminator. The LLM only
proposes group membership and writes the group description; the compiler
validates every proposal (members exist, no reuse, size bounds, no id
collisions) and synthesizes the union schema and routing table
deterministically, so a poor proposal can be rejected but never mis-routed.
Dispatch validates arguments against the \emph{member's} original schema. An
opt-in second phase merges overlapping capabilities \emph{across} servers
(two news providers) into capability facades whose actions name their source,
making vendor choice an explicit, documented decision.

\textbf{Macro inlining} adds hand-authored multi-step tools (typed steps over
existing tools) as first-class nodes, contributing data edges to the graph.

\textbf{Selection} replaces the served catalog with two meta-tools ---
\texttt{search\_tools} (BM25 retrieval over compiled definitions) and
\texttt{call\_tool} (validated dispatch by id) --- shrinking served context
to a few hundred tokens regardless of catalog size, at the price of a search
round trip.

\section{Serving Runtime}\label{sec:runtime}

The compiled artifact is served as a standard MCP endpoint. Every dispatch is
logged (arguments, result, latency, parent call for facade/macro/meta
fan-out), giving operators full-payload observability with live tail.

\textbf{Result compaction.} Results exceeding a token threshold are
summarized at serve time by a small model under instructions that preserve
identifiers and sourcing metadata (record ids, permalinks, timestamps), with
a deterministic structural fallback (bounded arrays, capped strings) when the
model is unavailable, and an asynchronous mode above a second threshold that
serves a structural cut immediately while caching the model-compacted version
for repeat calls. Error results and selection meta-results are never
compacted.

\textbf{Profile-guided re-optimization.} Recurring runtime errors are
clustered by (tool, normalized error signature); each new cluster is analyzed
by an LLM against the tool definition \emph{as agents currently see it}, and
the proposal is constrained to a typed action: amend the description (a
persistent overlay) or block the tool. Operators apply fixes with one action;
the applied overlay survives recompilation. The loop was validated by
reproducing, automatically, two fixes we had first made by hand (a missing
numeric constraint; a chronically failing tool).

\section{Evaluation}\label{sec:eval}

\textbf{Setup.} 137 tools across six live servers: a filesystem reference
server, the public Hugging Face and GitHub MCP servers, a deliberately naive
machine-generated National Weather Service server (one tool per GET
operation, verbatim descriptions), and two REST APIs (Open-Meteo,
Hacker News search) composed through our OpenAPI frontend --- chosen to
create genuine cross-server overlap in weather and article search. 60 tasks
across six categories (single-tool, ambiguous-selection, cross-server
chains, needle-in-haystack, no-tool-exists, and \emph{overlap-selection}:
tasks answerable by multiple sources, including sources that \emph{cannot}
answer, e.g.\ the US-only NWS for Paris weather), 3 trials each, two agent
models (\texttt{claude-sonnet-4-6}, \texttt{claude-haiku-4-5}; 1{,}800
trials total), answers judged by a stronger model against task-specific
rubrics. All conditions traverse the same gateway hop.

\begin{table}[h]
\centering
\begin{tabular}{lcccc}
\toprule
& \multicolumn{2}{c}{\texttt{sonnet-4-6}} & \multicolumn{2}{c}{\texttt{haiku-4-5}} \\
Condition & Success & \$/task & Success & \$/task \\
\midrule
Raw federation (137 tools) & 96.7\% & \$0.260 & \textbf{90.6\%} & \$0.090 \\
Compiled (rewrite+selection) & 96.1\% & \$0.064 & 77.2\% & \$0.027 \\
Compiled + consolidation & 96.1\% & \textbf{\$0.052} & 81.7\% & \textbf{\$0.022} \\
Consolidation, no selection & 94.4\% & \$0.151 & 80.6\% & \$0.107 \\
Cross-server consolidation & 96.1\% & \$0.054 & 81.7\% & \$0.026 \\
\bottomrule
\end{tabular}
\caption{Results over 60 tasks $\times$ 3 trials per condition per model.
Sonnet CIs overlap fully between raw and the compiled conditions.}
\end{table}

Four observations. First, for the stronger agent, \emph{compilation is
approximately free accuracy at one-fifth the cost}: all compiled conditions
are statistically indistinguishable from the raw federation, and compiled
consolidation wins both axes. Second, \emph{compilation repairs cross-source
arbitration}: on overlap-selection tasks the raw federation scored 90.0\%
with the stronger agent while every compiled condition scored 100\% ---
pool-aware disambiguation rewriting outperforms the naive federation exactly
where multiple vendors overlap. Third, and centrally, \emph{the benefit is
capability-gated}: the weaker agent reads the raw 25{,}000-token catalog
competently (90.6\%) but pays roughly nine points for compiled indirection
(its errored-call rate rises from 5.3\% to $\sim$24\%), and its
overlap-selection result inverts. Surface compression trades context for
indirection, and the price of indirection falls with agent capability ---
approaching zero for frontier models. Deployments should select surfaces per
agent model, which the toggle architecture makes a configuration choice.
Fourth, \emph{cross-server consolidation is safe but not yet differentiated}:
it ties compiled consolidation on both models; disambiguation rewriting
appears to carry the arbitration load that explicit capability merging
targets. A 30-task pilot preceding this study had suggested a 100\% ceiling
for the no-selection condition; at doubled task count that figure regressed
to 94.4\%, a reminder to distrust small-$n$ ceilings.

\textbf{Threats to validity.} Small task count and trial count (wide
intervals); judge circularity (a Claude model judging a Claude agent,
mitigated by rubric-based judging of final answers only and by human
spot-checks); live-service drift between condition snapshots (mitigated by
compiling all conditions from a single catalog snapshot); single agent model
(the larger study adds model diversity).

\section{Discussion}

The compiler framing pays off in governance as much as performance. Because
every transformation is an overlay over immutable originals, the system can
show its work: consoles render exact diffs, facades display per-action
provenance, and operator overrides compose with automated passes rather than
fighting them. The same properties make LLM participation safe: models
propose (rewrites, groupings, specification drafts, error fixes) while
deterministic validators decide, and every acceptance is reviewable and
reversible.

The runtime completes a loop that static compilation cannot: logged calls
reveal which tools agents actually use and how they fail; failures become
typed fix proposals; applied fixes recompile the surface. We believe this
observe--optimize--recompile loop, rather than any single pass, is the durable
contribution.

\section{Future Work}

Immediate directions: learned (log-driven) dead-tool elimination
and macro \emph{mining} from recurring call chains; schema-level rewriting
(folding parameter guidance into field descriptions); and compilation
targets beyond MCP.

\section*{Availability}

The engine is open source (Apache-2.0) at
\url{https://github.com/toolc-dev/toolc} and on npm as \texttt{@toolc/*}; the
hosted service is at \url{https://toolc.dev}. Benchmark tasks, transcripts,
and reports ship with the repository.

\end{document}
